\documentclass[usenames,dvipsnames,10pt]{beamer}
\usepackage[utf8]{inputenc}
\usepackage[T1]{fontenc}
\usepackage[german]{babel}
\usepackage{tikz}

%%%----------%%%----------%%%----------%%%----------%%%
\usepackage{listings}
\lstset{
    language=C,
    basicstyle=\ttfamily\small,
    aboveskip={1.0\baselineskip},
    belowskip={1.0\baselineskip},
    columns=fixed,
    extendedchars=true,
    breaklines=true,
    tabsize=4,
    prebreak=\raisebox{0ex}[0ex][0ex]{\ensuremath{\hookleftarrow}},
    frame=lines,
    showtabs=false,
    showspaces=false,
    showstringspaces=false,
    keywordstyle=\color[rgb]{0.627,0.126,0.941},
    commentstyle=\color[rgb]{0.133,0.545,0.133},
    stringstyle=\color[rgb]{01,0,0},
    numbers=left,
    numberstyle=\small,
    stepnumber=1,
    numbersep=10pt,
    captionpos=t,
    escapeinside={\%*}{*}
}

%%%----------%%%----------%%%----------%%%----------%%%
\title{Robotics - Are you going to be replaced?}
\date[ISPN ’80]{TE - \today}
\author{Eric Antosch}

\usetheme{hawsx}
\setbeamertemplate{section in toc}{%
  
      \begin{tikzpicture}
        
        \fill[color=UmUBlue] (1,1) circle (0.2cm);
        \node[anchor=center] (xlab) at (1,1) {\color{white}\inserttocsectionnumber};
        \node[anchor=west] (ylab) at (1.2,0.97) {\color{UmUBlue}\sffamily\textsf{\inserttocsection}};
      \end{tikzpicture}
  }


\begin{document}
    
    \begin{frame}
        \titlepage
    \end{frame}    
    \begin{frame}{Table of Contents}
        \tableofcontents
    \end{frame}
    \section{What is Robotics?}
    \framepic{Idee.jpg}{
        %\framefill
        \textcolor{white}{What is Robotics?}
        \vskip 0.5cm
    }
    \begin{frame}{Definition}
        \begin{block}{Definition}
            
        \end{block}
    \end{frame}
    \section{A brief history}
    
    \framepic{Idee.jpg}{
        %\framefill
        \textcolor{white}{A brief history}
        \vskip 0.5cm
    }
    \begin{frame}{A brief history}
        \begin{block}{Important facts}
            \begin{itemize}
                \item was Funktionen sind und wie man sie definiert,
                \pause
                \item was Headerdateien sind, welche es gibt und wie man eigene erstellt,
                \pause
                \item wo der Unterschied zwischen Call-By-Value und Call-By-Reference ist
                \pause
                \item und was Rekursion ist und wie man diese erstellt.
            \end{itemize}
        \end{block}
    \end{frame}

    \section{Overview}
    \subsection{Types}
    \framepic{Idee.jpg}{
        %\framefill
        \textcolor{white}{Overview: Types}
        \vskip 0.5cm
    }
    \subsection{Uses}
    \framepic{Idee.jpg}{
	%\framefill
    \textcolor{white}{Overview: Uses}
    \vskip 0.5cm
    }

    \begin{frame}{Overview}{Types}
        \begin{block}{Was ist eine Funktion?}
            Eine Funktion ist ein Block von Code, der durch einen Aufruf ausgeführt werden kann. Dabei besteht jede Funktion aus einem 
            Funktionskopf und einem Funktionskörper. Im Funktionskopf wird der Name, der Typ des Rückgabewerts und die Parameter festgelegt, während
            im Funktionskörper der eigentliche Code vorgefunden werden kann.
        \end{block}
    \end{frame}
    \begin{frame}{Overview}{Uses}
        \begin{block}{Prototypen}
            Es ist ratsam, alle Funktionen in dem eigenen Code direkt unter den Includes und den Präprozessordirektiven mithilfe von Prototypen
            zu beschreiben. Der Prototyp ist im Prinzip nur der Funktionskopf ohne Funktionskörper.
        \end{block}
    \end{frame}
    \subsection{Technical}
    \begin{frame}{Overview}{Technical}
        \begin{block}{Was sind Headerdateien?}
            Headerdateien sind der zweite Dateientyp neben *.c-Dateien, die ansich
            jedoch keinen C-Code enthalten. Sie enthalten lediglich die Informationen über tatsächliche
            C-Dateien, sodass sie in den eigenen Code eingebunden werden können.
        \end{block}
    \end{frame}

    \subsection{Economy}
    \begin{frame}{Overview}{Economy}
        \begin{block}{Call-By-Value}
            Der Aufruf einer Funktion mit Parametern, die Variablen
            darstellen. Die Werte werden in der Funktion zur Verarbeitung 
            von Anweisungen und Informationen genutzt.
        \end{block}
    \end{frame}

    \framepic{Idee.jpg}{
        %\framefill
        \textcolor{white}{Ethical Implications}
        \vskip 0.5cm
    }
    \section{Ethical Implications}
    \begin{frame}{Ethical Implications}{Relationships}
        \begin{block}{Was ist Rekursion?}
            Als eine rekursive Funktion bezeichnet man eine Funktion, die sich in ihrer Ausführung selbst aufruft.
            Man kann solche Funktionen gut mit rekursiven Folgen vergleichen. Das nächste Folgenglied ist hierbei abhängig von 
            den vorherigen Ergebnissen der Folge. Ein Anfang wird dabei meist als gegeben angenommen.
        \end{block}
    \end{frame}
    \begin{frame}{Ethical Implications}{Society}
        \begin{block}{Stack}
            Zu beachten bei dem Erstellen von rekursiven Funktionen ist, dass ein sogenannter Stack existiert. Ruft eine Funktion sich selbst auf,
            dann wird zuerst die zuletzt aufgerufene Funktion bis zum Schluss ausgeführt. Nach dem Ende dieser Funktion, springt das Programm zurück 
            in die Funktion, die die gerade beendete Funktion aufgerufen hat. Rekursion ist also nicht einfach nur eine andere Form von Iteration.
        \end{block}
    \end{frame}
    \begin{frame}{Ethical Implications}{Responsibility}
        
    \end{frame}


    \section{Article}

    \framepic{Idee.jpg}{
        %\framefill
        \textcolor{white}{Fragen?}
        \vskip 0.5cm
    }


    \section{Discussion}

    \begin{frame}{Discussion}
        \begin{exampleblock}{Key Questions}
            \begin{itemize}
                \item About what part of the future of robotics are you most excited about?
                \item Would you consider buying a robotic pet if it were indifferentiable from a normal one?
                \item In what part of society do robots belong? Does that social standing change with appearance or skill set?
                \item Is it possible for humans to become obsolete because of robotics?
                \item Are you going to be replaced?
            \end{itemize}
        \end{exampleblock}
    \end{frame}
    
    
    \section{Conclusion, Questions, Feedback}
    
    
    \framepic{Idee.jpg}{
        %\framefill
        \textcolor{white}{Conclusion}
        \vskip 0.5cm
    }

    \framepic{Idee.jpg}{
        %\framefill
        \textcolor{white}{Questions?}
        \vskip 0.5cm
    }

    \framepic{Idee.jpg}{
        %\framefill
        \textcolor{white}{Feedback}
        \vskip 0.5cm
    }

    \framepic{Idee.jpg}{
        %\framefill
        \textcolor{white}{Thank you!}
        \vskip 0.5cm
    }
    

\end{document}